\documentclass[12pt]{article}
\usepackage[utf8]{inputenc}
\usepackage[T1]{fontenc}
\usepackage{hyperref}
\title{Central Okinawan Language Resources for K-12 Curriculum Development}
\author{Compiled by Claude Sonnet 5.5 (claude-sonnet-5-5) for Andr\'as Kornai}
\date{}
\begin{document}
\maketitle

\section{Introduction}
Central Okinawan (Uchinaaguchi, ISO 639-3 ryu, Glottocode cent2126) is a Japonic language of Okinawa Island. It is endangered. More materials exist for it than for the other Ryukyuan languages, but the search results were mostly summaries and I could not check the details of each item.

\section{Primary Educational Resources}
\begin{itemize}
\item \emph{Basic Okinawan: From Conversation to Grammar}, by Rumiko Shinzato (with Shoichi Iwasaki), published by the University of Hawai\textquoteleft i Press. It is described as the first Okinawan textbook for English speakers, with ten units, and it does not assume knowledge of Japanese. Listings disagree on the publication date (2022 versus 2024). \url{https://scholarspace.manoa.hawaii.edu/items/019bfc00-a9f0-4561-b80d-4ba25037e974}
\item ``Rikka, Uchinaa-nkai! Okinawan Language Textbook for Beginners'' (2017), a beginner PDF shared through a blog post. Its provenance and copyright status are unclear. \url{https://liuchiuan.com/2017/09/13/rikka-uchinaa-nkai-okinawan-language-textbook-for-beginners-2017}
\item ``Let's Go To Okinawa'': an introduction to Uchinaguchi on the University of Hawai\textquoteleft i Okinawa pages. \url{https://manoa.hawaii.edu/okinawa/wordpress/?p=567}
\end{itemize}

\section{Linguistic and Reference Materials}
\begin{itemize}
\item University of Hawai\textquoteleft i ScholarSpace item on Uchinaguchi teaching materials. \url{https://scholarspace.manoa.hawaii.edu/handle/10125/5073}
\item Okinawa Prefecture overview (PDF) of the common kana orthography for the five Ryukyuan areas. \url{https://www.pref.okinawa.lg.jp/_res/projects/default_project/_page_/001/009/625/01_gaiyou.pdf}
\item Uruma City explanatory PDF on shimakutuba. I did not verify its contents beyond the search summary. \url{https://www.city.uruma.lg.jp/documents/10237/kaisetusimaku.pdf}
\item Tokyo University of Foreign Studies repository article (PDF) on Ryukyuan language. I did not verify its contents beyond the search summary. \url{https://tufs.repo.nii.ac.jp/record/2901/files/jo000407.pdf}
\item Dictionaries mentioned in search summaries but not individually linked: \emph{Okinawago jiten} (the National Institute for Japanese Language Okinawan dictionary), Sakihara's \emph{Okinawan-English Wordbook}, and the Yomitan village dictionary of about 15,000 entries.
\end{itemize}

\section{Religious and Community Resources}
\begin{itemize}
\item Ryukyu Shimpo article on the Uchinaguchi class of the Okinawa Association of America, a diaspora community programme. \url{https://english.ryukyushimpo.jp/2011/06/30/1551/}
\item Ryukyu Shimpo tag page on shimakutuba news. \url{https://ryukyushimpo.jp/tag/しまくとぅば}
\end{itemize}

\section{Digital and Audio Resources}
\begin{itemize}
\item The prefecture produces an animated series, ``Icharibachoo-ritsu Akijabi Elementary School,'' to promote shimakutuba. I saw it only in a summary and have no verified URL for it.
\end{itemize}

\section{Curriculum Development Context}
A news report in the search results says Okinawa Prefecture plans to set up an expert committee in fiscal 2026 and to prepare teaching materials and curriculum guidelines by fiscal 2027. The report is at \url{https://ryukyushimpo.jp/news/entry-1499643.html}. The prefecture's Shimakutuba Promotion Center keeps a human-resource bank that can supply instructors for school classes and teacher training. A 2026 volume by Nakamoto Ken and Murakami Rori, published by Sanseido, discusses language culture and school education in the Ryukyu Islands. A 2014 paper noted that learning materials for self-study were still insufficient. The Rethinking Schools essay ``Uchinaaguchi: The Language of My Heart'' covers the history of the language's exclusion from schools. \url{https://rethinkingschools.org/articles/uchinaaguchi-the-language-of-my-heart}
\end{document}